\begin{lemma}
For every positive integer $D$, there is a finite-vertex, finite-edge
$3$-uniform $2$-simple hypergraph $\mathcal F$ of maximum degree at most $D$
and a distinguished edge $f$ such that
\[
b_{L(\mathcal H),3D}(e)\le b_{L(\mathcal F),3D}(f)
\]
for every $3$-uniform $2$-simple hypergraph $\mathcal H$ with finitely many
edge indices, maximum degree at most $D$, and distinguished edge $e$.
The competitors' vertex sets need not be finite.
\end{lemma}
\begin{proof}
We construct a finite set of representatives, without maximizing over an
unbounded class of hypergraphs. Set $M=1+3(D-1)$ and $L=3M$.
Retain from a competitor only $e$ and its line-graph neighbors $N$.
For each of the three vertices of $e$, at most $D-1$ other edge indices
contain that vertex. Every member of $N$ contains at least one such vertex,
so $|N|\le3(D-1)$. Thus the retained edge-index set has size $m\le M$.
Restrict the vertex set to the union of the retained edges, of size
$\ell\le3m\le L$. Retained edges still have size three. Their pairwise
intersection sizes are unchanged, hence remain at most two for distinct
indices. At each retained vertex, its new degree counts a subset of its
old incident edge indices and is therefore at most $D$. These are exactly
the conditions in \noderef{HypergraphClass}.

The distinguished edge has exactly the same neighbors before and after
restriction. Intersections of two retained edges have not changed, so
adjacency between any two of these neighbors has not changed either, by
\noderef{LineGraphOfHypergraph}. The vertex degree in the line graph and
the independent pairs and triples in its neighborhood, as defined in
\noderef{IndependentPairCount} and \noderef{IndependentTripleCount}, are
consequently unchanged. The three terms in \noderef{LocalBParameter} show that the local
$b$-value is unchanged.

Relabel these $m$ edge indices by $\operatorname{Fin}(m)$ and these $\ell$
vertices by $\operatorname{Fin}(\ell)$. Bijections preserve incidences,
intersection cardinalities, degrees, and the three local counts just
described. In particular they preserve $2$-simplicity, not just
$3$-simplicity. Consider all choices of $0\le m\le M$, $0\le\ell\le L$,
all functions from $\operatorname{Fin}(m)$ to subsets of
$\operatorname{Fin}(\ell)$ satisfying the three class conditions, and all
distinguished indices in $\operatorname{Fin}(m)$. This is a finite set:
there are finitely many sizes, finitely many subsets of each finite vertex
set, finitely many such functions, and finitely many distinguished indices.
It is nonempty: take one three-element edge on three vertices, whose
vertex degrees are one and whose distinct-edge intersection condition is
vacuous. This is admissible since $D\ge1$, $M\ge1$, and $L\ge3$.

Apply mathlib's finite-domain maximum principle \texttt{Finite.exists\_max} to the real local
$b$-value on this finite nonempty set. Its maximizing element is a genuine
admissible hypergraph on two finite sets. Every competitor has the same
value as an element of this set by the restriction and relabeling above,
so its value is bounded by this maximum. No unused edge indices are padded
by duplicate hyperedges; the variable finite sizes preserve $2$-simplicity
throughout.
\end{proof}
